\section{Related Work}
\label{sec:related}

\paragraph{Withdrawal and consent on mobile.} MowChecker~\cite{du2024mowchecker} is the closest work. It statically
checks whether a withdrawal choice control- or data-flow-dominates third-party collection, that is, whether the toggle
is wired to the collection at all. Ours starts where a correct wiring
ends: data collected and queued \emph{before} the toggle, carried by persisted work across lifecycle events, and the
enforcement point that still sees it. The two are complementary. Studies of consent notices and explicit-consent
violations~\cite{nguyen2021sharefirst,nguyen2022freely,kollnig2021faitaccompli,matte2020cookie} and of tracking
ecosystems~\cite{razaghpanah2018trackers,kollnig2022goodbye} measure whether apps ask and whether they send before asking.
Policy-code consistency work~\cite{slavin2016policy,andow2019policylint,andow2020policheck,zimmeck2019maps} compares stated
policies with data flows; our oracle judges submissions against a stated contract in time.

\paragraph{Permissions and asynchronous use.} Android's permission model and its pitfalls are well
studied~\cite{felt2011demystified,felt2012attention,au2012pscout}. Aper~\cite{wang2022aper} finds dangerous APIs used
asynchronously after permission checks, a crash- and misuse-oriented cousin of our check-to-commit race. Apps also circumvent permissions through side channels~\cite{reardon2019fiftyways},
so we restrict our claims to cooperative apps and mediated paths.

\paragraph{Data management and information flow.} TaintDroid~\cite{enck2010taintdroid} tracks data to the network at
runtime. AppFence~\cite{hornyack2011appfence} blocks or shadows sensitive flows. Static analyzers such as
FlowDroid~\cite{arzt2014flowdroid}, IccTA~\cite{li2015iccta}, Amandroid~\cite{wei2014amandroid} and
DroidSafe~\cite{gordon2015droidsafe} find source-to-sink paths (see~\cite{li2017staticsurvey}). CleanOS~\cite{tang2012cleanos}
evicts idle sensitive data and Pebbles~\cite{spahn2014pebbles} recovers logical data objects. Weir~\cite{nadkarni2016weir} brings decentralized information-flow control~\cite{myers1997difc,efstathopoulos2005asbestos,zeldovich2006histar,krohn2007flume}
to Android, and FlowFence~\cite{fernandes2016flowfence} confines sensitive computation in IoT frameworks. Aurasium and
Boxify~\cite{xu2012aurasium,backes2015boxify} interpose on app behavior without platform changes, and
DroidForce~\cite{rasthofer2014droidforce} and AppGuard~\cite{backes2013appguard} enforce data-centric policies on Android
apps through inline reference monitors. Usage control~\cite{park2004ucon} re-evaluates a decision while use continues;
\sys's late check is such an ongoing decision, taken at every bounded write. Purpose-based access
control~\cite{byun2008purpose} ties access to the purpose of use, and Thoth~\cite{elnikety2016thoth} attaches policies to
data in server-side retrieval systems. \sys borrows the label idea, but its labels carry an \emph{epoch} that a withdrawal
invalidates, persist with \emph{job inputs} across reboot and force-stop, and feed a purpose- and pending-data-sensitive
decision.

\paragraph{Deletion, revocation and concurrency.} Secure deletion~\cite{reardon2013deletion}, formal deletion
definitions~\cite{garg2020deletion} and GDPR-oriented storage benchmarks~\cite{shastri2020gdprbench} address erasing data,
not stopping its later use by already-queued work. Our soundness argument is mutual exclusion: a policy
check and each bounded write are serialized against withdrawals, in the spirit of TOCTOU defenses~\cite{bishop1996races}. Leases~\cite{gray1989leases} bound how stale a cached grant can be, but do
not close the race on one device. Wake-lock studies~\cite{liu2016wakelock} examine background-work misuse for energy; we
study it for privacy contracts.

\paragraph{SDK isolation and downstream policy.} The SDK Runtime~\cite{android_sdkruntime} isolates SDKs but supplies no source epochs; Riverbed~\cite{wang2019riverbed} enforces data policies in attested server runtimes, beyond the device.
